\section{Formulación MILP}
\label{sec:formulacion_milp}

El problema de minimizar el número de S-boxes activas se formula mediante
Programación Lineal Entera Mixta. El modelo representa bit a bit dos
ejecuciones de Keccak reducido, denominadas izquierda y derecha, e impone
de manera exacta las transformaciones de cada ronda.

A diferencia de una formulación diferencial truncada, el modelo conserva
los valores absolutos de ambos estados. Las diferencias se calculan mediante
restricciones XOR y la propagación a través de $\chi$ se obtiene evaluando
la función booleana en las dos ejecuciones.

Esta representación incrementa el número de variables, pero permite
reconstruir parejas concretas de estados y verificar independientemente
las soluciones obtenidas por el solver.


\subsection{Índices y variables}
\label{subsec:variables_milp}

Los índices utilizados son

\[
r \in \{0,\ldots,R-1\},
\qquad
x,y \in \{0,\ldots,4\},
\qquad
k \in \{0,\ldots,z-1\},
\]

donde $r$ identifica la ronda, $(x,y)$ identifica una palabra del estado
y $k$ representa una posición dentro de la palabra.

Las dos ejecuciones se identifican mediante

\[
\sigma \in \{L,R\}.
\]

Para cada ejecución se introducen las variables binarias

\[
A_{r,x,y,k}^{\sigma} \in \{0,1\},
\]

que representan el estado al inicio de la ronda $r$.

Los estados intermedios se representan mediante

\[
T_{r,x,y,k}^{\sigma},
\qquad
R_{r,x,y,k}^{\sigma},
\qquad
B_{r,x,y,k}^{\sigma},
\qquad
C_{r,x,y,k}^{\sigma},
\]

donde:

\begin{itemize}
    \item $T^\sigma_r$ es la salida de $\theta$;
    \item $R^\sigma_r$ es la salida de $\rho$;
    \item $B^\sigma_r$ es la salida de $\pi$ y la entrada de $\chi$;
    \item $C^\sigma_r$ es la salida de $\chi$.
\end{itemize}

También se introducen variables de diferencia

\[
\Delta A_{r,x,y,k} \in \{0,1\},
\qquad
\Delta B_{r,x,y,k} \in \{0,1\},
\]

y variables de actividad

\[
a_{r,y,k} \in \{0,1\}.
\]


\subsection{Linealización de la operación XOR}
\label{subsec:xor_milp}

Sean $u$, $v$ y $d$ variables binarias relacionadas mediante

\[
d = u \oplus v.
\]

La operación XOR se representa exactamente introduciendo una variable
auxiliar binaria $q$ y la igualdad

\[
u+v=d+2q.
\label{eq:xor_lineal}
\]

Esta restricción reproduce las cuatro combinaciones posibles:

\[
\begin{array}{c|c|c}
u & v & d\\
\hline
0 & 0 & 0\\
0 & 1 & 1\\
1 & 0 & 1\\
1 & 1 & 0
\end{array}
\]

La misma construcción se utiliza en todas las operaciones XOR binarias
del modelo. Cuando es necesario calcular el XOR de más de dos bits, la
operación se descompone en una secuencia de XOR binarios.

Las diferencias entre las dos ejecuciones se definen mediante

\[
\Delta A_{r,x,y,k}
=
A_{r,x,y,k}^{L}
\oplus
A_{r,x,y,k}^{R},
\]

y

\[
\Delta B_{r,x,y,k}
=
B_{r,x,y,k}^{L}
\oplus
B_{r,x,y,k}^{R}.
\]


\subsection{Representación de las capas lineales}
\label{subsec:capas_lineales_milp}

Las transformaciones $\theta$, $\rho$ y $\pi$ se incorporan de manera
independiente para cada ejecución $\sigma$.


\subsubsection{Transformación \texorpdfstring{$\theta$}{theta}}

Para cada columna se calcula la paridad

\[
P_{r,x,k}^{\sigma}
=
\bigoplus_{y=0}^{4}
A_{r,x,y,k}^{\sigma}.
\]

En la implementación, este XOR de cinco bits se construye mediante
operaciones XOR binarias sucesivas.

El término de difusión se define como

\[
D_{r,x,k}^{\sigma}
=
P_{r,(x-1)\bmod 5,k}^{\sigma}
\oplus
P_{r,(x+1)\bmod 5,(k-1)\bmod z}^{\sigma}.
\]

La salida de $\theta$ satisface

\[
T_{r,x,y,k}^{\sigma}
=
A_{r,x,y,k}^{\sigma}
\oplus
D_{r,x,k}^{\sigma}.
\]

Todas estas relaciones se linealizan utilizando la
Ecuación~\ref{eq:xor_lineal}.


\subsubsection{Transformaciones \texorpdfstring{$\rho$}{rho} y
\texorpdfstring{$\pi$}{pi}}

Las transformaciones $\rho$ y $\pi$ son permutaciones de bits. Por tanto,
se representan mediante igualdades directas y no requieren variables
auxiliares de producto.

La transformación $\rho$ se impone mediante

\[
R_{r,x,y,k}^{\sigma}
=
T_{r,x,y,(k-r[x,y])\bmod z}^{\sigma}.
\]

La transformación $\pi$ se representa como

\[
B_{r,y,(2x+3y)\bmod 5,k}^{\sigma}
=
R_{r,x,y,k}^{\sigma}.
\]

Estas restricciones conservan el valor de cada bit y modifican únicamente
su posición dentro del estado.


\subsection{Representación exacta de la capa \texorpdfstring{$\chi$}{chi}}
\label{subsec:chi_milp}

Para cada ejecución $\sigma$, la transformación $\chi$ se define por

\[
C_{r,x,y,k}^{\sigma}
=
B_{r,x,y,k}^{\sigma}
\oplus
\left[
\left(
1-B_{r,(x+1)\bmod 5,y,k}^{\sigma}
\right)
B_{r,(x+2)\bmod 5,y,k}^{\sigma}
\right].
\]

La expresión contiene el producto de dos variables binarias. Para
linealizarlo se introduce

\[
s_{r,x,y,k}^{\sigma} \in \{0,1\},
\]

con

\[
s_{r,x,y,k}^{\sigma}
=
\left(
1-B_{r,(x+1)\bmod 5,y,k}^{\sigma}
\right)
B_{r,(x+2)\bmod 5,y,k}^{\sigma}.
\]

La conjunción se representa mediante las restricciones

\[
s_{r,x,y,k}^{\sigma}
\leq
1-B_{r,(x+1)\bmod 5,y,k}^{\sigma},
\]

\[
s_{r,x,y,k}^{\sigma}
\leq
B_{r,(x+2)\bmod 5,y,k}^{\sigma},
\]

\[
s_{r,x,y,k}^{\sigma}
\geq
B_{r,(x+2)\bmod 5,y,k}^{\sigma}
-
B_{r,(x+1)\bmod 5,y,k}^{\sigma}.
\]

Como todas las variables son binarias, estas desigualdades fuerzan
exactamente el valor del producto.

La salida de $\chi$ se obtiene mediante

\[
C_{r,x,y,k}^{\sigma}
=
B_{r,x,y,k}^{\sigma}
\oplus
s_{r,x,y,k}^{\sigma},
\]

que se linealiza nuevamente mediante la
Ecuación~\ref{eq:xor_lineal}.

La formulación se aplica por separado a las ejecuciones izquierda y
derecha. La diferencia de salida surge, por tanto, de dos evaluaciones
exactas de la función booleana $\chi$.


\subsection{Representación de la capa \texorpdfstring{$\iota$}{iota}}
\label{subsec:iota_milp}

La constante de ronda es conocida y no constituye una variable de decisión.

Para la palabra $(0,0)$, si el bit correspondiente de la constante es cero,
se impone

\[
A_{r+1,0,0,k}^{\sigma}
=
C_{r,0,0,k}^{\sigma}.
\]

Si el bit de la constante es uno, se utiliza

\[
A_{r+1,0,0,k}^{\sigma}
=
1-C_{r,0,0,k}^{\sigma}.
\]

Para las demás posiciones se imponen igualdades directas:

\[
A_{r+1,x,y,k}^{\sigma}
=
C_{r,x,y,k}^{\sigma},
\qquad
(x,y)\neq(0,0).
\]

Aunque la misma constante se cancela al calcular la diferencia entre
ejecuciones, su inclusión permite representar y reconstruir correctamente
los estados absolutos.


\subsection{Variables de actividad}
\label{subsec:variables_actividad_milp}

La variable

\[
a_{r,y,k} \in \{0,1\}
\]

indica si la aplicación local de $\chi$ situada en $(y,k)$ se encuentra
activa durante la ronda $r$.

Esta variable debe valer uno cuando al menos uno de los cinco bits de
entrada presenta una diferencia no nula. La relación OR se representa
mediante

\[
a_{r,y,k}
\geq
\Delta B_{r,x,y,k},
\qquad
x\in\{0,\ldots,4\},
\]

y

\[
a_{r,y,k}
\leq
\sum_{x=0}^{4}
\Delta B_{r,x,y,k}.
\]

Si los cinco bits de diferencia son cero, la segunda restricción fuerza

\[
a_{r,y,k}=0.
\]

Si al menos uno de los bits es uno, alguna de las cotas inferiores fuerza

\[
a_{r,y,k}=1.
\]

El número de S-boxes activas en la ronda $r$ es

\[
m_r
=
\sum_{y=0}^{4}
\sum_{k=0}^{z-1}
a_{r,y,k}.
\]


\subsection{Condición de diferencia inicial no nula}
\label{subsec:diferencia_inicial_milp}

Sin una restricción adicional, el modelo admitiría como solución trivial
dos estados iniciales idénticos, cuya actividad total sería cero.

Para excluir esta solución se impone

\[
\sum_{x=0}^{4}
\sum_{y=0}^{4}
\sum_{k=0}^{z-1}
\Delta A_{0,x,y,k}
\geq
1.
\]

Esta condición garantiza que

\[
A_0^L \neq A_0^R.
\]


\subsection{Función objetivo}
\label{subsec:funcion_objetivo_milp}

La actividad total durante $R$ rondas se define como

\[
N_{\mathrm{act}}
=
\sum_{r=0}^{R-1}
\sum_{y=0}^{4}
\sum_{k=0}^{z-1}
a_{r,y,k}.
\]

La función objetivo es

\[
\min
\sum_{r=0}^{R-1}
\sum_{y=0}^{4}
\sum_{k=0}^{z-1}
a_{r,y,k}.
\label{eq:funcion_objetivo}
\]

Cuando el solver demuestra la optimalidad global, el valor de la
Ecuación~\ref{eq:funcion_objetivo} corresponde a

\[
N_{\mathrm{act}}^{\min}(z,R).
\]

Cuando se obtiene únicamente una solución factible, su actividad constituye
una cota superior para el mínimo global.


\subsection{Restricciones auxiliares}
\label{subsec:restricciones_auxiliares}

Además del problema de minimización, la implementación permite incorporar
restricciones auxiliares para certificar cotas y validar testigos.


\subsubsection{Cota sobre la actividad total}

Para comprobar si existe una trayectoria con actividad no mayor que un
valor $U$, se añade

\[
N_{\mathrm{act}} \leq U.
\]

Si el solver demuestra que este problema es infactible, se concluye que no
existe una trayectoria válida con actividad total menor o igual que $U$.


\subsubsection{Actividad exacta por ronda}

El número de S-boxes activas en una ronda concreta puede fijarse mediante

\[
m_r
=
\sum_{y=0}^{4}
\sum_{k=0}^{z-1}
a_{r,y,k}
=
\widehat{m}_r,
\]

donde $\widehat{m}_r$ es el valor impuesto.

Esta restricción permite estudiar distribuciones como

\[
2+2
\]

o

\[
2+2+c.
\]


\subsubsection{Fijación de soportes activos}

Sea

\[
S_r
\subseteq
\{0,\ldots,4\}
\times
\{0,\ldots,z-1\}
\]

un soporte activo conocido. Puede imponerse

\[
a_{r,y,k}=1
\qquad
\text{si }(y,k)\in S_r,
\]

y

\[
a_{r,y,k}=0
\qquad
\text{si }(y,k)\notin S_r.
\]

Esta operación se utiliza para comprobar con el modelo MILP los soportes
obtenidos mediante búsqueda o reconstrucción.


\subsubsection{Fijación de estados concretos}

Los estados iniciales de un testigo pueden fijarse mediante

\[
A_{0,x,y,k}^{L}
=
\widehat{A}_{0,x,y,k}^{L},
\]

\[
A_{0,x,y,k}^{R}
=
\widehat{A}_{0,x,y,k}^{R}.
\]

Con ambos estados fijados, el modelo verifica que las transformaciones, las
diferencias y la actividad reportada sean consistentes con todas las
restricciones.


\subsection{Interpretación de los resultados del solver}
\label{subsec:interpretacion_solver}

La interpretación del resultado depende del estado de terminación del
solver:

\begin{description}
    \item[Óptimo global.] Se encontró una solución factible y se demostró
    que no existe otra con menor valor objetivo.

    \item[Solución factible.] Se encontró un testigo válido, pero no se
    demostró la optimalidad. Su valor constituye una cota superior.

    \item[Infactible.] Se demostró que ninguna solución satisface las
    restricciones impuestas. Este resultado puede utilizarse para elevar
    una cota inferior.

    \item[Límite de tiempo.] La terminación por tiempo no demuestra por sí
    sola optimalidad, infactibilidad ni ausencia de soluciones.
\end{description}

Esta distinción es fundamental en los experimentos de tres rondas. Los
testigos con actividades 13 y 14 demuestran cotas superiores globales,
pero su optimalidad solo fue establecida dentro de la familia restringida

\[
2+2+c.
\]